% ======================= RESUMEN =======================
\section*{Resumen}

\textbf{Problema y objetivo.} El Sistema Integral de Monitoreo Ambiental (SIMA) de Nuevo León registra con frecuencia días con PM2.5 elevado en la Zona Metropolitana de Monterrey (ZMM). Nos preguntamos cuánto mejora la meteorología observada del mismo día la identificación de los días-estación con PM2.5 promedio de 24\,h mayor de 25\ugm{} (límite de la NOM-025-SSA1-2021, usado como referencia fija), frente a un modelo que solo conoce la estación y la temporada.

\textbf{Datos y preparación.} Usamos los datos horarios de 13 estaciones entre 2021 y 2025. La limpieza por reglas físicas y de consistencia con la red invalidó el 0.7\,\% de los valores. \textbf{La variable respuesta no se imputa:} el PM2.5 diario se calcula solo con al menos 18 horas observadas. La meteorología horaria se imputó con métodos de imputeTS elegidos por enmascaramiento en 2021–2023; los días con hasta cuatro variables faltantes se completaron con la mediana de la red ajustada por estación, validada contra MICE (errores casi iguales; con MICE la conclusión no cambia). La base final tiene 17\,828 días-estación; los días sin PM2.5 suficiente se concentran en NO, NE2 y SUR y en 2024–2025.

\textbf{Técnicas.} Componentes principales (técnica de interdependencia) para describir la estructura de la meteorología, y dos técnicas de dependencia en paralelo para comparar M0 (estación y temporada) con M1 (M0 más ocho variables meteorológicas): el análisis discriminante lineal (LDA), fijado en el protocolo, y la regresión logística, nuestra referencia para interpretar las variables. Entrenamos con 2021–2023 y evaluamos 2024–2025, años que ya se habían usado en etapas previas (Sección~\ref{sec:supuestos}). La incertidumbre se estimó con un bootstrap pareado por bloques de 14 días (el tamaño del bloque, no un horizonte de pronóstico), también estratificado por temporada: es una validación temporal útil, no una confirmación independiente.

\textbf{Resultados.} La meteorología del día mejora el AUC entre 0.071 (regresión logística) y 0.079 (LDA) sobre estación y temporada, en la clasificación retrospectiva de 2024–2025 (LDA: de 0.713 a 0.792, IC 95\,\% [0.054, 0.104]). M1 detecta el 74\,\% de los días con superación (M0: 56\,\%) y el 72\,\% de los episodios, con una especificidad algo menor (0.68 frente a 0.70). Frente a M0, la mejora es mayor de marzo a mayo (0.151), no se distingue de cero en lluvias, se mantiene entre 0.071 y 0.082 en las variantes analizadas y sube a 0.090 si la logística admite relaciones no lineales. Con signo estable en los dos clasificadores, la superación se asocia con viento débil, presión más baja de lo normal para el sitio y, al ajustar por estación y temporada, mayor temperatura. La persistencia (PM2.5 del día anterior) clasifica mejor que la meteorología del día (AUC 0.872 frente a 0.792); sobre ella, la meteorología agrega 0.016, sobre todo de noviembre a febrero.

\textbf{Qué sí y qué no se puede concluir.} En el periodo evaluado, la meteorología observada mejora la clasificación retrospectiva frente a estación y temporada; no se respalda un pronóstico, no se identifican causas y las probabilidades no están calibradas. El LDA no cumple normalidad ni igualdad de covarianzas y los errores siguen autocorrelacionados: las medidas tomadas acotan la incertidumbre y evalúan el desempeño, pero no convierten el análisis en un modelo temporal ni en una alerta operativa.

\textbf{Recomendaciones para SIMA.} Evaluar el modelo con meteorología pronosticada y con la persistencia como base, junto con su área de pronósticos; elegir el corte según el uso (0.18 para vigilancia, con sensibilidad de 0.84 en la prueba, o 0.26 para equilibrar errores); incorporar dos reglas de control de calidad (radiación de madrugada y viento contra la red), y revisar el PM2.5 de García (NO2) entre 2023 y 2025 y la radiación de NO2 y NTE.

\clearpage
